\section{The two zero-ternary frame exceptions}
\label{sec:zero-ternary-frames}
Put
\[
 z_2(G)=\max_{L\lhd G}\dim_{\F_2}L/[L,G]L^2,
 \qquad
 \mu_T(M)=\max_{B\leq M\ {\mathrm{invariant}}}\dim_{\F_2}B/[B,T].
\]
The second quantity is the maximum dimension of a trivial homogeneous
section of the one actual module $M$.

\begin{lemma}[A fixed-frame bound]\label{lem:frame-rank}
If $1\to E\to U\to T\to1$ has elementary binary kernel
$E\leq M$ as a $T$-module, then
$z_2(U)\leq\mu_T(M)+z_2(T)$. Moreover
\[
 z_2(U)=\max_{N\lhd U}\operatorname{rank}\Omega_1(Z(U/N))_2.
\]
\end{lemma}
\begin{proof}
For $K\lhd U$, let $L$ be its image in $T$, put
$N=[K,U]K^2$ and $I=K\cap E$. The image of $N$ is exactly
$[L,T]L^2$, so there is an exact sequence
\[
 0\longrightarrow I/(I\cap N)\longrightarrow K/N
       \longrightarrow L/[L,T]L^2\longrightarrow0.
\]
Its first term is a trivial quotient of the actual submodule $I$
of $M$, proving the bound. Each $K/[K,U]K^2$ is central elementary
binary in its corresponding quotient of $U$. Conversely the preimage
of any central elementary binary subgroup in $U/N$ has relative
head surjecting onto that subgroup. Taking maxima proves the identity.
\end{proof}

\begin{lemma}[Noncentral socle types in the two frames]
\label{lem:small-power-socle}
Suppose $U$ has exact local $A_4$ or $S_4$ on $s\in\{8,16\}$
blocks and zero ternary linear kernel. Let $M=(\F_2^2)^s$,
$E=U\cap M$ and $T=U/E$, its actual faithful block top.
Every nontrivial binary simple type in the socle of every quotient
has Schur demand at most $5s/8$, and every odd-primary type has
demand at most $\lfloor s/3\rfloor$. Thus every quotient has a
faithful tuple of at most $5s/8$ characters if $z_2(U)\leq5s/8$.
\end{lemma}
\begin{proof}
In $Q=U/N$, the whole binary socle $A$ is a semisimple $T$-module:
the normal binary $\bar E$ acts trivially on each simple type.
Filter $A$ by $A\cap\bar E$. For a type $X$ of binary dimension
$d$ and Schur dimension $e$, its multiplicities satisfy
\[
 m_E\leq\lfloor2s/d\rfloor,\qquad
 m_T\leq\lfloor s/(2d)\rfloor.
\]
The second bound uses the whole elementary outer section of $T$.
These multiplicities add, since the types on the two sides may agree.
For $d\geq4$, their sum is at most $5s/8$. For $d=3,e=3$,
dividing the sum by three and rounding up also suffices.
For $d=3,e=1$, the image is $C_7$; its normal kernel is transitive,
since its orbit number divides both seven and $s$. Restricted to
that kernel, $X^{m_E}$ is trivial of dimension $3m_E$.
Lemma~\ref{lem:block-homogeneous} bounds this dimension by $s$,
giving $m_E\leq\lfloor s/3\rfloor$ and
$m_T\leq\lfloor s/6\rfloor$.

For $d=2,e=2$ the demand is at most
$\lceil(s+s/4)/2\rceil=5s/8$. If $d=2,e=1$, the image is
$C_3$ and its normal kernel $K$ is transitive. At a block
stabilizer, restriction to $K$ either keeps a nontrivial local
$C_3$ or $S_3$, or kills the local $C_3$ in the $A_4$ case.
In the first case take a Sylow $3$-subgroup $P$ of $K$. At a
$P$-fixed block it is Sylow in the stabilizer and its local image
is $C_3$, with no fixed vector. Every nonfixed orbit contributes
at most two invariant dimensions. Exactness of odd-order invariants
on sections gives $2m_E\leq2\lfloor s/3\rfloor$.
Adding $m_T\leq s/4$ suffices at both counts.
In the second case the induced frame restricted to $K$ is two
copies of the honest permutation module $\F_2^s$. Its transitive
Sylow $2$-subgroup gives the scalar submodule capacity
$\binom a{\lfloor a/2\rfloor}=3s/8$ for $s=2^a\in\{8,16\}$.
Filtering a trivial section through the two copies gives
$m_E\leq3s/8$, and again $m_E+m_T\leq5s/8$.

For odd $p$, the entire $p$-socle meets $\bar E$ trivially and
embeds in a quotient of $T$, so its rank is at most $\lfloor s/p\rfloor$.
The only remaining binary type is the trivial type, whose maximum
rank over the normal quotients is exactly $z_2(U)$ by
Lemma~\ref{lem:frame-rank}. Apply the whole-socle criterion.
\end{proof}

The following finite inputs concern tops and induced linear frames,
and do not enumerate all extensions $U$ above a frame. They are
supplied with literal permutation, module and kernel witnesses.
\begin{lemma}[Finite frame inputs]\label{lem:finite-zero-frames}
The zero-ternary frame certificates have the following properties.
\begin{enumerate}[label=(\roman*)]
\item Every nonbinary transitive top $T$ of degree eight satisfies
$z_2(T)\leq2$. Every frame induced from an onto map of a point
stabilizer to $C_3$ or $S_3$ satisfies $\mu_T(M)\leq2$.
\item For each degree-sixteen top having four four-point blocks,
exact local $A_4$ or $S_4$, and zero internal ternary kernel,
every induced $C_3/S_3$ frame admits a Sylow $3$-subgroup $P$
and a $2$-element $g\in N_T(P)$ with
$z_2(T)+\dim(M^P)^g\leq10$.
\item Every primitive degree-sixteen top and every eligible induced
frame admit an odd subgroup $A$ with $z_2(T)+\dim M^A\leq10$.
\item The primitive degree-eight actions are the three affine groups
$2^3:C_7$, $2^3:(C_7:C_3)$, $\operatorname{AGL}_3(2)$ and
$\operatorname{PSL}_2(7)$, $\operatorname{PGL}_2(7)$, $A_8,S_8$.
Their normal subgroup orders are respectively
\[
 (1,8,56),\ (1,8,56,168),\ (1,8,1344),\ (1,168),\
 (1,168,336),\ (1,20160),\ (1,20160,40320).
\]
Except for the first group, a Sylow $3$-subgroup has orbit lengths
$1,1,3,3$ in each actual action.
\end{enumerate}
\end{lemma}
\begin{proof}
The degree-eight top alphabet consists of the seventeen imprimitive
classes constructed by the four-pair and two-four-block Goursat and
lift certificates, and the seven primitive actions of the stated
primitive classification. All normal subgroups are retained for the
relative-rank check. All onto point-stabilizer maps are recovered by
their literal kernels; for a fixed kernel, automorphisms of $C_3$
are realized by simultaneous Frobenius and automorphisms of $S_3$
are inner. This yields all thirty frames on twenty eligible tops.

For every one of these frames the supplied binary invariant-submodule
list contains zero and is closed under adding every cyclic submodule.
The latter are constructed from all vector orbits under the actual
generator matrices. Any omitted invariant submodule would contain
an included one maximal under containment; adjoining the cyclic
submodule of a vector outside it contradicts maximality. This proves
completeness, and direct commutator ranks give the stated maximum two.

For (ii), the complete induced-frame, kernel and complement construction
gives twenty-nine four-by-four top actions, including their nonsplit
lifts. Each has exactly one eligible stabilizer kernel. The supplied
$P,g$ witnesses give the inequality directly; no invariant-submodule
enumeration in dimension 32 is required. For (iii), all twenty-two
primitive degree-sixteen entries are checked, giving nine eligible
stabilizer kernels on seven tops and the stated odd-subgroup witnesses.
The primitive completeness premise is
\cite{RoneyDougal2005Primitive,PrimGrp344}; parts (i), (ii) use the
separate constructive completeness arguments just described. The
normal orders and actual Sylow orbits in (iv) are direct checks on
the same seven primitive representatives. Thus these are finite
mathematical premises with independently checkable completeness
certificates, rather than an inference from their cardinalities.
\end{proof}

\begin{theorem}[Complete zero-ternary eight and sixteen frames]
\label{thm:zero-ternary-small-frames}
Every exact-local $A_4/S_4$ action with zero ternary kernel and
eight or sixteen blocks has a faithful character cover on every
normal quotient of length five or ten, respectively. Their entire
original-normal menus satisfy Corollary~\ref{cor:character-forward}.
\end{theorem}
\begin{proof}
By Lemmas~\ref{lem:frame-rank} and \ref{lem:small-power-socle},
it suffices to bound $\mu_T(M)+z_2(T)$ by five or ten.
At degree eight Lemma~\ref{lem:finite-zero-frames}(i) gives four.

We record two tools for degree sixteen. Odd-order invariants are
exact, so for an odd subgroup $P$,
$\mu_T(M)\leq\mu_{N_T(P)}(M^P)$. If $g$ is a $2$-element
normalizing $P$, the binary cyclic module $M^P\downarrow_{\langle g\rangle}$
is a sum of modules over a chain ring. The head and socle have
equal dimension, and submodules cannot have larger socle. Hence
\[
 \mu_T(M)\leq\dim(M^P)^g.
\]
Also, if $P$ is Sylow three in a transitive top of power-of-two
degree, a $P$-fixed block contributes zero to $M^P$: the local
onto $C_3/S_3$ map has nontrivial $C_3$ image there. Each nonfixed
$P$-orbit contributes at most two invariant dimensions.

Suppose first $T$ preserves eight pairs, with kernel $K$ and top
$S\leq S_8$. If $K\ne1$, it projects onto every pair. The
local map from a point stabilizer kills its intersection with $K$,
because $C_3,S_3$ have no nontrivial normal $2$-subgroup. It
descends onto the stabilizer of a pair in $S$, giving one of the
degree-eight frames $M_0$. Induction in stages and the length-two
filtration of the binary two-point permutation module give a
filtration of $M$ with two inflated copies of $M_0$. Thus
$\mu_T(M)\leq2\mu_S(M_0)\leq4$.
Furthermore Lemma~\ref{lem:frame-rank} on $K\leq\F_2^8$ gives
$z_2(T)\leq3+z_2(S)\leq5$, using the scalar induced capacity
three in degree eight. The sum is at most nine.
If $K=1$, then $T\cong S$ and $z_2(T)\leq2$. A Sylow three
has at most two nonfixed orbits on the eight pairs, each lifting
to two orbits on the sixteen points, since an odd pair stabilizer
fixes its pair pointwise. Hence $\mu_T(M)\leq8$, giving ten.

Otherwise a minimal nontrivial block of $T$ has size four or eight,
or $T$ is primitive. For size four let $F$ be its internal binary
translation kernel, let $D$ be its internal ternary kernel and let
$S\leq S_4$ be the top. If $D\ne1$, retained-$F$ quotients
have character cover two by Theorem~\ref{thm:four-nonzero-core}.
For quotients killing $F$, the central elementary binary socle
injects after removing the odd $D$ into a quotient of
$C_2\wr S\leq S_8$. Therefore $z_2(T)\leq4$.
If $S=A_4$ or $S_4$, a Sylow three has at most three nonfixed
orbits on the sixteen points: over its three cycled four-blocks
the moving local triples give one orbit of length nine and their
fixed points one orbit of length three; the fourth block contributes
at most one further orbit. Thus $\mu_T(M)\leq6$.
If $S$ is binary, let $P$ be Sylow three of the full block kernel.
It is also Sylow in $T$, and Frattini gives a normalizer transitive
on the four blocks. Each block contributes the same dimension
$d\leq2$ to $M^P$. The normalizer modulo $P$ is binary, so the
local $d$-space has a series of trivial factors. Induction gives
$d$ copies of the four-point permutation module in a filtration
of $M^P$. Their scalar capacity is two, giving
$\mu_T(M)\leq2d\leq4$. If instead the internal $D=1$, use
Lemma~\ref{lem:finite-zero-frames}(ii) and the normalizer tool.

For a minimal block of size eight, the two-block kernel is a
Goursat subdirect product of one of the seven primitive locals in
Lemma~\ref{lem:finite-zero-frames}(iv). The first, $2^3:C_7$,
has no eligible point-stabilizer map. For each other local a Sylow
three has four nonfixed orbits on the sixteen points, so
$\mu_T(M)\leq8$. For $\operatorname{PGL}_2(7)$ or $S_8$,
the intersection of the kernel with the square of the simple
socle is semisimple, and the quotient embeds in $C_2\wr C_2$
of degree four; hence $z_2(T)\leq2$.
For the simple locals the block kernel is semisimple, so its
relative binary rank is zero. For the two eligible affine locals,
the displayed normal menus show that every intermediate Goursat
kernel is either a graph of the local group, the full product, or
contains both translations with odd or common simple quotient.
The translations are then a semisimple sum of nontrivial simple
modules, including all automorphism twists. They have no trivial
section, and the quotient has relative binary rank zero. Thus the
kernel has $z_2=0$ and $z_2(T)\leq1$. Every case gives a sum
at most ten. The primitive degree-sixteen case is
Lemma~\ref{lem:finite-zero-frames}(iii).

These cases exhaust the top: a primitive binary group has prime
degree, so no omitted binary primitive local occurs at sizes four
or eight. The original extension $U$ is arbitrary throughout.
Finally $5^{64\cdot5}<2^{24\cdot32-3}$ and
$5^{64\cdot10}<2^{24\cdot64-3}$ give the character-forward
margin at the two physical widths.
\end{proof}
